\section{Energy-selected survival for a family of cells}
\label{sec:tpc50-frame-selection}

The fixed-shape theorem below controls every nonzero standard cell.
It is useful to have a second statement that remains valid when a
partition contains boundary atoms with arbitrarily small energy.  The
result is elementary, but it separates low cell energy from poor mask
survival without selecting a favorable signed summand.

Let \(M:I\times J\to\{0,1\}\), where \(|J|=N\), and suppose every row
has at most \(\epsilon N\) forbidden entries.  Let
\(b_\lambda\in\C^J\), \(\lambda\in\Lambda\), be a finite cell family.
Fix \(B>0\) and \(A>0\), and assume
\begin{equation}
 |\Lambda|\le S,
 \qquad
 \|b_\lambda\|_\infty^2\le B^2,
 \qquad
 \sum_{\lambda\in\Lambda}\|b_\lambda\|_2^2
 \ge ANB^2.
 \label{eq:tpc50-frame-hypotheses}
\end{equation}
The last inequality is only an aggregate lower frame bound.  It does
not assert orthogonality of the cells or a lower bound for their signed
sum.

For \(0<\tau<1\), call a cell energetic if
\begin{equation}
 \|b_\lambda\|_2^2\ge\tau NB^2.
 \label{eq:tpc50-energetic-cell}
\end{equation}

\begin{theorem}[Energy selection without a cellwise lower bound]
\label{thm:tpc50-energy-selection}
Every energetic cell and every nonzero \(a\in\C^I\) satisfy
\begin{equation}
 \delta_w(M;a,b_\lambda)\ge1-\frac\epsilon\tau.
 \label{eq:tpc50-energetic-survival}
\end{equation}
Moreover, the total energy of the nonenergetic cells is bounded by
\begin{equation}
 \frac{\sum_{\lambda\text{ nonenergetic}}
             \|b_\lambda\|_2^2}
      {\sum_{\lambda}\|b_\lambda\|_2^2}
 \le\frac{S\tau}{A}.
 \label{eq:tpc50-low-cell-energy}
\end{equation}
Consequently, if
\(\epsilon=X^{-c+o(1)}\), \(S/A=X^{o(1)}\), and
\(\tau=\epsilon^{1/2}\), then every energetic cell has survival at least
\(1-X^{-c/2+o(1)}\), while the nonenergetic cells carry only
\(X^{-c/2+o(1)}\) of the aggregate cell energy.
\end{theorem}

\begin{proof}
For an energetic cell,
\begin{equation}
 \operatorname{Flat}(b_\lambda)
 =\frac{N\|b_\lambda\|_\infty^2}
       {\|b_\lambda\|_2^2}
 \le\frac1\tau.
 \label{eq:tpc50-energy-flatness}
\end{equation}
The degree--flatness inequality gives
\eqref{eq:tpc50-energetic-survival}.  Each nonenergetic cell has energy
less than \(\tau NB^2\); there are at most \(S\) of them.  Divide their
sum by the aggregate lower bound in
\eqref{eq:tpc50-frame-hypotheses} to obtain
\eqref{eq:tpc50-low-cell-energy}.  The last assertion is immediate.
\end{proof}

\begin{corollary}[Bounded-overlap partitions]
\label{cor:tpc50-partition-frame}
Suppose functions \(w_\lambda:J\to\C\) have overlap at most \(R\),
satisfy
\begin{equation}
 \sum_\lambda w_\lambda(j)=1
 \quad(j\in J),
 \qquad
 |w_\lambda(j)|\le C,
 \label{eq:tpc50-partition-unity}
\end{equation}
and let \(0\ne g\in\C^J\) obey
\begin{equation}
 \|g\|_2^2\ge cN\|g\|_\infty^2.
 \label{eq:tpc50-base-flatness}
\end{equation}
Then \(b_\lambda=gw_\lambda\) satisfies
\eqref{eq:tpc50-frame-hypotheses} with
\(S=|\Lambda|\), \(B=C\|g\|_\infty\), and
\(A=c/(RC^2)\) (or any smaller positive value).
\end{corollary}

\begin{proof}
At every \(j\), Cauchy--Schwarz and the overlap bound give
\begin{equation}
 1=\left|\sum_\lambda w_\lambda(j)\right|^2
 \le R\sum_\lambda|w_\lambda(j)|^2.
 \label{eq:tpc50-square-frame-pointwise}
\end{equation}
Multiply by \(|g_j|^2\) and sum over \(j\).  Thus

\[
 \sum_\lambda\|gw_\lambda\|_2^2
 \ge R^{-1}\|g\|_2^2
 \ge (c/R)N\|g\|_\infty^2
 =\frac{c}{RC^2}NB^2.
\]

The pointwise upper bound is immediate.
\end{proof}

TPC--18 uses a fixed bounded-overlap dyadic partition in the opened
\(d\)-variable \citep{WangTPC18}; fixed residue cells have overlap one.
Provided the base-flatness and aggregate lower-energy hypotheses in
\cref{cor:tpc50-partition-frame} are verified for the family at hand,
\cref{thm:tpc50-energy-selection} then protects its energy-significant
boundary atoms even if a uniform lower bound is not imposed on every
individual cutoff.  It is deliberately an aggregate cell-energy
statement.  Turning that aggregate into the squared norm of a signed
Mellin reassembly is a separate frame/no-cancellation problem.
